
\chapter{Neural Network Components: Optimizers}
\label{Optimizers}


\begin{figure}[!ht]
    \centering
    \includegraphics[width=\textwidth]{../../images/08_Optimizers_ From Gradients to Updates.png
    }
\end{figure}

\paragraph{}
We've already established the basic workings of an
end-to-end neural network. We've seen how tensors,
operations, graphs, layers and losses come together
to perform a sequence of mathematical operations
that produce the intended output.

However, measuring how far the prediction is from
the desired result is only part of the process. As
discussed in Chapter \ref{Losses}, the loss function
quantifies the prediction error. We still need a way
to minimize this error by updating the model
parameters. This is where optimizers come into play.

During initialization, the model parameters are
typically assigned values using an initialization
strategy. These initial values generally do not
produce satisfactory predictions. During training,
the optimizer uses the gradients calculated during
backpropagation to update the model parameters,
helping the network minimize its loss.

\paragraph{}
We've already seen that the model parameters
are initialized randomly. So at first it is
going to give very incorrect output. Now it
is optimizer which makes changes to the
model parameters based on the gradients
calculated during backward pass that pushes
model parameters to calculate output close to
desired target. Lets see how it works.

\section{Gradient Descent: GD}
Gradient descent is one of the fundamental
optimization algorithms used to train neural
networks. It updates model parameters in the
direction opposite to the gradient of the loss
function.

Compared with second-order methods such as
Newton-Raphson, standard gradient descent is
generally simpler to implement and requires less
memory because it does not need to compute and
store the full Hessian matrix. This makes it a
practical choice for large neural networks.

\[ W_\text{t+1} = W_\text{t} - \eta * \partial{W_t} \]

from the above example we can easily tell
what a optimizer suppose to do. For a
model parameter $W$ we're updating it's value
by subtracting partial difference times $\eta$
with original value. $\eta$ is learning rate.

\paragraph{Learning Rate}
The learning rate, denoted by $\eta$, controls the
step size of each parameter update.

If the learning rate is too large, the optimizer
may overshoot the minimum, oscillate around it,
or even diverge. If it is too small, convergence
may become unnecessarily slow.

Choosing an appropriate learning rate is therefore
critical to effective training. In practice, the
learning rate may also be adjusted during training
using a learning-rate schedule.

We're not going to work on a single sample.
So working with a large sample size. Based
on how we can tackle a large sample size
and update model parameter values there are
three distinct GD stratigies in place.

\begin{itemize}
    \item Batch Gradient Descent(BGD)
    \item Mini Batch Gradient Descent(MGD)
    \item Stocastic Gradicnt Descent(SGD)
\end{itemize}

Now, Lets take a look at each of them.

\subsection{Batch Gradient Descent: BGD}
The gradient is calculated using the entire
training dataset, and the parameters are
updated once after the full dataset has
been evaluated.

\paragraph{Pros:} The gradient represents
the entire dataset, giving a consistent
update direction for the current parameter
values.

\paragraph{Cons:} Each update can be
computationally expensive for large
datasets.

\subsection{Mini Batch Gradient Descent: MGD}
The training dataset is divided into smaller
batches. The model processes each batch and
updates its parameters after each batch.


\paragraph{Pros:} Balances computational
efficiency and gradient stability. It
also works well with vectorized CPU and
GPU operations.

\paragraph{Cons:} The gradient is an
estimate of the full-dataset gradient,
so updates can fluctuate between batches.
\subsection{Stocastic Gradient Descent: SGD}
In its strict definition, the gradient is
calculated using one randomly selected
training sample per update.


\paragraph{Pros:} Updates are inexpensive and
frequent. The noise in the gradient can
sometimes help the optimizer move away
from saddle regions or shallow local minima.

\paragraph{Cons:} Updates can fluctuate
considerably, making convergence less
smooth.

\section{Stocastic Gradient Descent}
\label{Optimizer::SGD}

I've briefly touched SGD in the previous
section but lets dig deeper into this
algorithm. So I've already established
how SGD works Mathematically. Let's see
how it is integrated with the framework.



Optimizers are selected through the registry and used through the
shared-owning \texttt{tf::optimizer::optimizer} handle. String keys replace the
old enum-based selection, and an unknown key is rejected when the handle or
\texttt{model.compile()} is called.

\begin{lstlisting}[
    caption={Selecting an optimizer through the public API},
    label={lst:optimizer public api}
]
auto sgd = tf::optimizer::optimizer(tf::optimizer::type::sgd);

tf::optimizer::config config;
config.set(std::string(tf::optimizer::parameter::learning_rate), 0.05);
auto configured = tf::optimizer::optimizer(tf::optimizer::type::sgd, config);
configured.setParameter(tf::optimizer::parameter::learning_rate, 0.02);

tf::model model({x}, dense_output);
model.compile(configured, tf::loss::type::squared_error);
model.compile("sgd", tf::loss::type::squared_error);
\end{lstlisting}

The handle owns the optimizer instance together with the compiled model, so the
optimizer stays alive while either one is in use. \texttt{setParameter()}
dispatches on the optimizer's canonical type, which also works for optimizers
created through a registry alias.

Internally, \texttt{Optimizer} is still an abstract base class, but its graph
and update buffers are implementation details. Layers and models never receive
an internal optimizer object; they call the public handle, and
\texttt{OptimizerRegistry} routes those calls to the registered implementation.

\begin{lstlisting}[
    caption={Internal optimizer contract},
    label={lst:Declear Optimizer class}
]
// file optimizer.hpp
class Optimizer {
protected:
  // Discards the previous update graph and its owned buffers.
  void resetOptimizerGraph();
  virtual void clearUpdateState() {}

public:
  virtual void createParameterUpdateGraph(tf::tensor param,
                                          tf::tensor &update_param,
                                          tf::tensor gradient_param) = 0;

  // Prepares this optimizer for a fresh training graph.
  void reset();
};
\end{lstlisting}

The \texttt{createParameterUpdateGraph()} subroutine is where this framework
differs from optimizers that work eagerly. Update operations are recorded in a
dedicated graph during training setup and executed after the forward and
backward passes. The graph is destroyed safely when the optimizer is replaced or
destroyed: API-owned tensors are spared first, and the concrete optimizer drops
any buffers that kept update tensors alive.

\begin{lstlisting}[
    caption={Defining SGD class},
    label={lst:SGD class}
]
// file sgd.cpp
SGD::SGD() : learning_rate(0.01) {}

void SGD::createParameterUpdateGraph(tf::tensor param, tf::tensor &update_param,
                                     tf::tensor gradient_param) {
  // RAII guard keeps GraphManager's stack balanced even on exceptions.
  OptimizerGraphScope scope(this->optimizer_graph);

  UpdateState &state = this->update_states.emplace_back();
  state.param = param;
  state.gradient = gradient_param;
  state.scaled_gradient.tf_create({param.getDimensions()[0],
                                   param.getDimensions()[1]},
                                  param.dt_type);

  state.scaled_gradient = gradient_param.scale(learning_rate);
  update_param = state.param.sub(state.scaled_gradient);
}
\end{lstlisting}

The RAII guard pushes the optimizer's internal graph onto the graph manager for
the duration of update recording and always balances the stack afterwards. Each
call owns its parameter, gradient, and scaled-gradient tensors in a typed update
state instead of a growing untyped buffer list. Calling \texttt{reset()} before a
new training run clears those states and builds a fresh update graph, so
repeated \texttt{compile()} and \texttt{fit()} calls do not reuse stale ops.

\paragraph{}
Although vanilla SGD is simple and computationally
efficient, it has some limitations. Its updates can
fluctuate significantly, and it may converge slowly
in regions where the loss landscape has poorly
conditioned curvature.

Several optimizer variants address these issues by
incorporating additional information into the
parameter updates. One of the simplest and most
important extensions is SGD with momentum, which
we will explore next.


\section{SGD with momentum}

